\begin{figure}[!htbp]
\centering
\begin{tikzpicture}
\begin{axis}[
    ybar,
    width=0.86\linewidth, height=0.52\linewidth,
    ymin=2.6, ymax=4.0,
    ylabel={Participant-category mean necessity},
    symbolic x coords={Collection,Sharing},
    xtick=data,
    legend style={at={(0.5,-0.22)},anchor=north,legend columns=3,draw=none,
                  font=\footnotesize},
    bar width=22pt,
    nodes near coords, point meta=y,
    nodes near coords={\pgfmathprintnumber[fixed,precision=2]{\pgfplotspointmeta}},
    every node near coord/.append style={font=\scriptsize,color=ngSlate},
    axis line style={ngSlate}, tick style={ngSlate},
    grid=major, grid style={ngSlate!18}
]
\addplot+[fill=ngRed!75,  draw=ngRed]
   coordinates {(Collection,3.14)(Sharing,3.27)};
\addplot+[fill=ngOrange!75, draw=ngOrange]
   coordinates {(Collection,3.18)(Sharing,3.47)};
\addplot+[fill=ngBlue!75,  draw=ngBlue]
   coordinates {(Collection,3.05)(Sharing,3.36)};
\legend{Survey: Personal, Survey: Maybe, Survey: Not personal}
\end{axis}
\end{tikzpicture}
\caption{The Necessity Gap, participant-linked. Within the \emph{same}
participants, classifying a category as Personal does not produce lower
necessity ratings: cluster-bootstrap differences are 0.09 [\textendash 0.27,
0.45] for collection and \textendash 0.09 [\textendash 0.44, 0.25] for sharing.}
\label{fig:rq5-necessity-gap}
\end{figure}